\title{FlashAttention-3: Fast and Accurate Attention with Asynchrony and Low-precision}

\begin{abstract}
The attention mechanism, integral to the pervasive Transformer framework, serves as the principal performance constraint in scaling large language models and extending to longer input sequences. \textsc{FlashAttention} introduced a method for accelerating attention on GPUs by optimizing memory read and write operations. Nonetheless, this approach has not fully leveraged the novel features introduced in the latest GPU hardware; for example, \textsc{FlashAttention-2} utilizes only 35\% of the H100 GPU’s computational capacity. To address this, we propose three key optimizations for improving attention performance on Hopper GPUs: (1) making use of the asynchronous operation between Tensor Cores and TMA through warp-specialized scheduling to concurrently manage computation and data transfer; (2) alternating the execution of block-level matrix multiplication and softmax calculations; and (3) implementing block quantization with incoherent processing to exploit FP8 low-precision arithmetic available on modern hardware. Our approach, termed \textsc{FlashAttention-3}, demonstrates 1.5-2.0$\times$ higher throughput on H100 GPUs, achieving up to 740 TFLOPs/s (representing 75\% utilization) in FP16 mode, and nearly 1.2 PFLOPs/s in FP8. Furthermore, we show that FP8 \textsc{FlashAttention-3} offers 2.6$\times$ reduced numerical error compared to conventional FP8 attention baselines.
\end{abstract}

\section{Introduction}
\label{sec:intro}

Within the Transformer framework~\citep{vaswani2017attention}, the attention module emerges as the chief computational constraint, as the calculation of self-attention among query and key representations scales quadratically with sequence length. Enabling attention to manage substantially longer contexts promises to facilitate a range of enhanced capabilities—including supporting modeling and inference across multiple extensive documents~\citep{guo2021longt5,shaham2022scrolls,peng2023yarn} and handling files within vast code repositories~\citep{roziere2023code, li2023starcoder}. It may also pave the way for accommodating novel data modalities such as detailed image inputs~\citep{chen2022scaling}, as well as audio~\citep{gulati2020conformer} and video sequences~\citep{ho2022video}, in addition to unlocking new usage scenarios like capturing prolonged user interactions~\citep{sun2019bert4rec} and supporting agent operations spanning extended decision horizons~\citep{yao2022react}. This backdrop has spurred extensive research efforts to accelerate attention in settings with long contexts, involving various approaches such as employing approximate algorithms~\citep{katharopoulos2020transformers,choromanski2020rethinking, tay2020efficient}, advancing software-level optimizations~\citep{rabe2021self,dao2022flashattention,kwon2023efficient}, and exploring fundamentally different architectural alternatives~\citep{peng2023rwkv,sun2023retentive,gu2023mamba}.

This paper extends the advancements made by \citet{dao2022flashattention}, who crafted precise attention algorithms informed by the architecture and execution details of GPUs. Their work, as presented in \citep{dao2022flashattention}, introduced \textsc{FlashAttention}, which employed an innovative tiling approach for attention parallelization. This technique integrated all attention-related computations into a unified GPU kernel, avoiding extra transfers to and from global memory and thereby significantly increasing computational efficiency. Later, \citet{dao2023flashattention2} enhanced the algorithm, releasing \textsc{FlashAttention-2}, which supports parallelization along the sequence length axis and processes key and value matrix blocks in the forward pass's inner loop. These modifications bolster load balancing and maximize GPU occupancy. Despite these improvements, our experiments reveal that \textsc{FlashAttention-2} underutilizes modern GPUs compared to state-of-the-art matrix-multiplication (GEMM) kernels—showing, for example, only 35\% utilization on the Hopper H100 GPU, as opposed to GEMM's typical 80-90\%. Some of this inefficiency stems from practical considerations; for instance, the algorithm relies on Ampere-era instructions when targeting Tensor Cores instead of the Hopper generation's more suitable operations. Previous research, including ThunkerKitten~\citep{spector2024thunder} and cuDNN 9~\citep{cudnn9}, demonstrates that leveraging Hopper-oriented instructions and employing tiling-based design principles can both accelerate attention computations and simplify code structure.

At a more foundational level, the algorithm behind \textsc{FlashAttention-2} operates within a streamlined synchronous framework, without intentionally incorporating asynchronous execution or low-precision arithmetic into its methodology. Asynchronous execution arises naturally from the deployment of specialized hardware components designed to optimize critical operations in machine learning tasks: matrix multiplications are handled by dedicated units like Tensor Cores, while separate units such as the Tensor Memory Accelerator (TMA) manage memory loading, distinct from general-purpose CUDA cores responsible for logical, integer, and floating-point operations. The adoption of low-precision formats—including FP8 in the Hopper architecture and FP4 in Blackwell, building on earlier usage of FP16 (in Pascal, 2017) and BF16 (in Ampere, 2020)—is an established method for substantially increasing throughput (by factors of two or four), without raising power consumption or silicon area requirements. We discuss how Hopper supports these advancements in more detail in \cref{subsec:hardware}. The central technical obstacle involves reengineering \textsc{FlashAttention-2} to harness these hardware developments: leveraging asynchrony means enabling computational overlap between matrix multiplications and softmax calculations—even though there are dependencies between them—while utilizing low-precision demands careful management to reduce quantization errors, particularly when addressing outlier features in large language models~\citep{dettmers2208llm, sun2024massive}.

To address these challenges, we introduce \textsc{FlashAttention-3}, which integrates and develops three novel strategies to further advance performance on modern GPU platforms:\footnote{Although our results are situated in the context of NVIDIA’s Hopper architecture, the presented algorithm extends to any GPU architecture that provides robust support for asynchronous execution and low-precision operations.}

\begin{enumerate}

\item \textbf{Producer-Consumer asynchrony:} We propose a warp-specialized form of software pipelining that leverages the decoupled execution capabilities of data transfers and Tensor Cores. By allocating data producers and data consumers to different warps, this scheme enhances the pipeline’s effectiveness in masking latencies associated with memory accesses and instruction issuance.

\item \textbf{Hiding softmax under asynchronous block-wise GEMMs:} Non-GEMM operations within the softmax calculation—such as floating point multiply-adds and exponentials—typically have lower throughput. We mask their cost by overlapping their execution with asynchronous WGMMA-based GEMMs. To achieve this, our implementation restructures the \textsc{FlashAttention-2} routine to resolve certain ordering dependencies between softmax and GEMM steps. As an illustration, in our two-stage variant, while a softmax operation is applied to a current scores matrix block, WGMMA proceeds asynchronously to process the subsequent block.

\item \textbf{Hardware-accelerated low-precision GEMM:} We refine the forward pass algorithm to utilize GEMMs performed on FP8 Tensor Cores, leading to a near twofold increase in observed TFLOPs/s. Since WGMMA stipulates specific memory layouts for FP32 accumulator and FP8 operand blocks, we resolve layout mismatches through block quantization and incoherent processing methods. These adjustments help maintain accuracy while exploiting the performance benefits of FP8 precision.

To empirically assess our approach, we evaluated \textsc{FlashAttention-3} on the H100 SXM5 GPU across various configurations, and our findings indicate the following: (1) The FP16 variant yields a 1.5-2.0$\times$ improvement in forward pass throughput relative to \textsc{FlashAttention-2} (achieving rates up to 740 TFLOPs/s), and a 1.5-1.75$\times$ gain in backward pass performance; (2) FP8 nearly reaches 1.2 PFLOPs/s; and (3) for tasks with extended sequence lengths, FP16 demonstrates superior efficiency, while FP8 shows comparable performance\footnote{Specifically, for head dimension 64, FP8 in \textsc{FlashAttention-3} leads, whereas for head dimensions 128 and 256, performance matches state-of-the-art implementations in the absence of causal masking, but trails when causal masking is involved.} with NVIDIA cuDNN's latest attention operation. Further, FP16 computations in \textsc{FlashAttention-3} maintain parity in numerical error with \textsc{FlashAttention-2}, and surpass conventional attention implementations due to preserving intermediate calculations (such as softmax rescaling) in FP32. Additionally, employing block quantization and incoherent processing in FP8 mode provides \textsc{FlashAttention-3} a 2.6$\times$ accuracy advantage over standard attention when per-tensor quantization is used, notably in scenarios with outlier features.

We have made \textsc{FlashAttention-3} available as open-source software under a permissive license\footnote{\textsc{FlashAttention-3}
  can be found at \texttt{https://github.com/Dao-AILab/flash-attention}}, with plans to integrate it into the PyTorch and Hugging Face ecosystems for maximum accessibility to the research and developer communities.

\section{Background: Multi-Head Attention and GPU Characteristics}
\label{sec:background}

\subsection{Multi-Head Attention}
\label{subsec:multi_head_attn}

Consider $\mathbf{Q}, \mathbf{K}, \mathbf{V} \in \mathbb{R}^{N \times d}$, which represent the query, key, and value matrices for a single attention head, where $N$ denotes the length of the input sequence and $d$ corresponds to the dimensionality per head. The attention mechanism produces the output $\mathbf{O}$ through the following steps:

\begin{equation*}
  \mathbf{S} = \alpha \mathbf{Q} \mathbf{K}^\top \in \mathbb{R}^{N \times N}, \quad \mathbf{P} = \mathrm{softmax}(\mathbf{S}) \in \mathbb{R}^{N \times N}, \quad \mathbf{O} = \mathbf{P}\mathbf{V} \in \mathbb{R}^{N \times d},
\end{equation*}

In this context, $\mathrm{softmax}$ operates across rows, and it is common to choose the scaling constant as $\alpha = 1/\sqrt{d}$. To avoid numerical issues arising from exponentiation, it is standard practice to subtract $\mathrm{rowmax}(\mathbf{S})$ from $\mathbf{S}$. In the case of multi-head attention (MHA), separate projection weights for queries, keys, and values are used for each head, allowing this procedure to be executed concurrently over all heads and batch elements, resulting in the complete output tensor.

Consider a scalar loss function $\phi$, and denote its gradient with respect to an argument as $\mathbf{d}(-) = \partial \phi / \partial (-)$. Provided with the output gradient $\mathbf{dO} \in \mathbb{R}^{N \times d}$, we utilize the chain rule to determine $\mathbf{dQ}$, $\mathbf{dK}$, and $\mathbf{dV}$ in the following manner:

\begin{align*}
  \mathbf{dV} &= \mathbf{P}^\top \mathbf{dO} \in \mathbb{R}^{N \times d} \\
  \mathbf{dP} &= \mathbf{dO} \mathbf{V}^\top \in \mathbb{R}^{N \times N} \\
  \mathbf{dS} &= \mathrm{dsoftmax} (\mathbf{dP}) \in \mathbb{R}^{N \times N} \\
  \mathbf{dQ} &= \alpha \mathbf{dS} \mathbf{K} \in \mathbb{R}^{N \times d} \\
  \mathbf{dK} &= \alpha \mathbf{dS}^\top \mathbf{Q} \in \mathbb{R}^{N \times d},
\end{align*}

Here, for a vector $s$ and its softmax $p = \mathrm{softmax}(s)$, the relationship $\mathbf{d}s = (\mathrm{diag}(p) - p p^\top)\mathbf{d}p$ holds; the operation $\mathrm{dsoftmax}(\mathbf{dP})$ corresponds to applying this transformation to each row individually. Moreover, these gradients for the backward computation of MHA are computed in parallel both over different heads and batches.

\subsection{GPU hardware characteristics and execution model}
\label{subsec:hardware}

We describe the aspects of the GPU's execution model relevant for \textsc{FlashAttention-3}, with a focus on the NVIDIA Hopper architecture as a concrete instantiation of this model.

\paragraph{Memory hierarchy:} The GPU features a hierarchical structure of memory locations, where storage size typically decreases as bandwidth increases (\cref{tab:gpu-hierarchy})\footnote{\citet{luo2024benchmarking} indicates that shared memory can deliver a bandwidth of 128 bytes per clock cycle for each SM; by considering 132 SMs and a boost clock frequency of 1830 MHz, we derive the aggregate bandwidth.}. At the base, global memory (GMEM), or HBM, provides off-chip DRAM available to all streaming multiprocessors (SMs). Data situated in GMEM is automatically placed in an on-chip L2 cache for faster access. Each SM further includes its own compact, highly banked on-chip shared memory (SMEM), which requires explicit management by programmers. The deepest and fastest tier in this hierarchy is the register file found inside every SM.

\paragraph{Thread hierarchy:} The organization of the GPU's programming framework relies on structured collections of execution units, referred to as threads. Progressing from the most granular to the broadest classification, this hierarchy consists of individual threads, warps (each consisting of 32 threads), warpgroups (comprised of 4 sequential warps), threadblocks (also known as cooperative thread arrays or CTAs), threadblock clusters (specific to Hopper), and ultimately, grids.

The two hierarchies exhibit a strong connection. Threads belonging to a particular CTA are executed concurrently on the same SM, while CTAs within a given cluster are jointly executed on a shared GPC. All threads in a CTA have direct access to SMEM, but each individual thread possesses a maximum of 256 dedicated registers (RMEM) exclusive to it.

\paragraph{Asynchrony and warp-specialization:}

GPUs function as throughput-oriented processors that depend on parallelism and asynchronous operations to mask both memory and computational delays. To facilitate asynchronous data transfers between GMEM and SMEM, the Hopper architecture incorporates a specialized hardware component known as the Tensor Memory Accelerator (TMA) \cite[\S7.29]{cuda}. Additionally, Hopper introduces an asynchronous Tensor Core—accessible through the warpgroup-level WGMMA instruction \cite[\S9.7.14]{ptx}—which, unlike its predecessor Ampere, is capable of directly obtaining operands from shared memory.

Asynchronous hardware features enable warp-specialized kernel design, wherein warps within a CTA are designated distinctly as either producers (focused solely on data transfers) or consumers (engaged strictly in computational tasks). This separation allows compilers to optimize instruction ordering more effectively \citep{warp-specialization-2011}. Furthermore, the Hopper architecture introduces the ability to redistribute registers dynamically among warpgroups via the \verb|setmaxnreg| directive \cite[\S9.7.17.1]{ptx}. As a result, warps performing matrix multiply-accumulate (MMA) operations can be allocated more RMEM, in contrast to warps only managing TMA, which requires just a single active thread.

\paragraph{Low-precision number formats:}
\label{sec:low-precision-gpu}
Contemporary GPUs come equipped with hardware modules tailored to accelerate calculations at lower numeric precisions. For instance, the WGMMA operation can exploit Hopper's FP8 Tensor Cores, achieving twice the per-SM throughput relative to computations performed in FP16 or BF16 formats.

Nonetheless, properly utilizing FP8 WGMMA requires awareness of the operand layout limitations. Consider a GEMM operation computing $A \times B^{\top}$, where $A$ is sized $M \times K$ and $B$ is $N \times K$. An operand is designated as \emph{mn-major} if its data is stored contiguously along the outer $M$ or $N$ dimensions, while it is labeled \emph{k-major} if the inner $K$-dimension is contiguous instead. In the case of FP16 WGMMA, operands located in SMEM can be provided in either mn-major or k-major layout. By contrast, FP8 WGMMA exclusively permits input operands in k-major arrangement. Additionally, scenarios—such as attention mechanisms—where there is a need to chain GEMMs within a single kernel can be hindered due to layout mismatches between FP32 accumulator outputs and the FP8 input requirement, making consecutive dependent FP8 WGMMA invocations problematic.

In the context of attention, these layout restrictions entail certain modifications to the design of an FP8 algorithm, which we describe in \cref{sec:algofp8}.

\subsection{Standard Attention and Flash Attention} In line with \citet{dao2022flashattention}, we use the term \textbf{standard attention} to refer to a GPU-based attention mechanism where the intermediate matrices $\mathbf{S}$ and $\mathbf{P}$ are explicitly stored in HBM. The central innovation of \textsc{FlashAttention} is the introduction of a localized softmax reduction, which mitigates the high cost associated with reading and writing these intermediate results by combining the attention process within a unified kernel. This local softmax reduction is implemented in lines \ref{code-ws:softmax_start}-\ref{code-ws:softmax_end} of the consumer mainloop in \cref{alg:flash3_wgmma_ws_only}, along with the necessary block-wise scaling of $\mathbf{O}$. A straightforward derivation confirming that this approach correctly computes $\mathbf{O}$ can be found in \cite[\S 2.3.1]{dao2023flashattention2}.

\section{FlashAttention-3: Algorithm}
\label{sec:algo}

In this section, we present the \textsc{FlashAttention-3} algorithm. To streamline the discussion, we concentrate on the forward pass; the procedure for the backward pass is outlined in~\cref{sec:algo_ws_bwd}. Initially, we demonstrate the integration of warp-specialization and the use of a circular SMEM buffer into the foundational \textsc{FlashAttention-2} algorithm. Next, we detail how the asynchronous characteristics of WGMMA enable the construction of a two-stage pipeline that overlaps GEMM and softmax computations. Finally, we outline the adjustments required for FP8, addressing both layout compliance and enhanced numerical accuracy through block quantization and handling incoherent computation.

\subsection{Producer-Consumer asynchrony through warp-specialization and pingpong scheduling}
\label{sec:algo_ws}

\paragraph{Warp-specialization}
Similar to \textsc{FlashAttention-2}, the forward computation in \textsc{FlashAttention-3} can be efficiently parallelized over the batch dimension, number of heads, and the length of the query sequence. Consequently, it is sufficient to describe the algorithm from a CTA (Cooperative Thread Array) perspective, focusing on the processing of a specific query matrix tile, $\mathbf{Q}_i$, to generate its associated output tile, $\mathbf{O}_i$. For clarity, we initially present the warp-specialization approach utilizing a circular SMEM buffer, omitting any discussion of overlapping the GEMM and softmax operations. Letting $d$ denote the head dimension and $N$ the length of the sequence, we select a query block size $B_r$ such that the matrix $\mathbf{Q}$ is divided into $T_r = \lceil \frac{N}{B_r} \rceil$ blocks, denoted $\mathbf{Q}_1, ..., \mathbf{Q}_{T_r}$.

\begin{algorithm}[H]
    \caption{\small\label{alg:flash3_wgmma_ws_only}\textsc{FlashAttention-3} forward pass \textbf{excluding} intra-consumer overlap -- CTA perspective}
    \begin{algorithmic}[1]
\REQUIRE Matrices $\mathbf{Q}_i \in \mathbb{R}^{B_r \times d}$ and $\mathbf{K}, \mathbf{V} \in \mathbb{R}^{N \times d}$ in HBM, key block dimension $B_c$ with $T_c = \lceil \frac{N}{B_c} \rceil$.
\STATE Set up pipeline structure to oversee barrier coordination utilizing $s$-stage circular SMEM buffers.
\IF {inside producer warpgroup}
\STATE Free a fixed amount of registers as planned.
\STATE Begin transfer of $\mathbf{Q}_i$ from HBM to shared memory.
\STATE After transfer, confirm to consumer that $\mathbf{Q}_i$ is available.
\FOR{$0 \le j < T_c$}
    \STATE Wait until the consumer finishes with the $(j\,\%\,s)$ buffer stage.
    \STATE Initiate loading of $\mathbf{K}_j, \mathbf{V}_j$ from HBM into the shared memory at stage $(j\,\%\,s)$.
    \STATE Upon completion, notify consumers that $\mathbf{K}_j$ and $\mathbf{V}_j$ are loaded.
\ENDFOR
\ELSE \STATE Reassign the previously freed registers according to the count of consumer warps.
\STATE Locally, set $\mathbf{O}_i = (0) \in \mathbb{R}^{B_r \times d}$ with $\ell_i = 0$, $m_i = -\infty$ in $\mathbb{R}^{B_r}$.
\STATE Await availability of $\mathbf{Q}_i$ in shared memory.
\FOR{$0 \le j < T_c$}
\STATE Wait for $\mathbf{K}_j$ data to arrive in shared memory.
\STATE Carry out matrix multiply $\mathbf{S}_i^{(j)} = \mathbf{Q}_i \mathbf{K}_j^T$ (SS-GEMM). Commit and wait. \label{alg:ws_only_gemm1}
\STATE Save $m_i^{\mathrm{old}} = m_i$; update $m_i = \mathrm{max}(m_i^{\mathrm{old}}, \mathrm{rowmax}(\mathbf{S}_i^{(j)}))$. \label{code-ws:softmax_start}
\STATE Calculate $\widetilde{\mathbf{P}}_i^{(j)} = \mathrm{exp}(\mathbf{S}_i^{(j)} - m_i)$ and update $\ell_i = \mathrm{exp}(m_i^{\mathrm{old}} - m_i) \ell_i + \mathrm{rowsum}(\widetilde{\mathbf{P}}_i^{(j)})$. \label{code-ws:softmax_end}
\STATE Wait until $\mathbf{V}_j$ resides in shared memory.
\STATE Update $\mathbf{O}_i = \mathrm{diag}(\mathrm{exp}(m_i^{\mathrm{old}} - m_i))^{-1} \mathbf{O}_i + \widetilde{\mathbf{P}}_i^{(j)} \mathbf{V}_j$ (RS-GEMM). Commit and wait. \label{alg:ws_only_gemm2}
\STATE Mark the $(j\,\%\,s)$ buffer stage as free for producer use.
\ENDFOR
\STATE Compute final outputs $\mathbf{O}_i = \mathrm{diag}(\ell_i)^{-1} \mathbf{O}_i$ and $L_i = m_i + \log(\ell_i)$.
\STATE Save $\mathbf{O}_i$ and $L_i$ to HBM as the $i$th output block in $\mathbf{O}$ and $L$.
\ENDIF
\end{algorithmic}
\end{algorithm}

In our implementation of \cref{alg:flash3_wgmma_ws_only} targeting Hopper, we utilize \verb|setmaxnreg| to manage (de)allocations and employ TMA for fetching $\mathbf{Q}_i$ along with $\{ \mathbf{K}_j, \mathbf{V}_j \}_{0 \leq j < T_c}$. GEMM operations within the consumer mainloop are performed via WGMMA, with the SS or RS designation reflecting whether the leading operand resides in shared memory or the register file, respectively. To understand the control flow of \cref{alg:flash3_wgmma_ws_only}, it is important to recognize that asynchronous TMA loads can be initiated without waiting for prior loads to finish. Additionally, within the producer mainloop, the absence of wait instructions during the initial $s$ iterations enables buffer priming.

\paragraph{Pingpong scheduling} The asynchronous execution provided by WGMMA and TMA, combined with warp specialization, enables the possibility to interleave the softmax calculation for one warpgroup while another warpgroup is executing GEMM. This strategy becomes compelling when considering that, on state-of-the-art hardware accelerators, non-matrix multiplication operations exhibit significantly lower throughput compared to matmul. For instance, the H100 SXM5 GPU is capable of delivering 989 TFLOPS for FP16 matrix multiplications, but achieves only 3.9 TFLOPS for specialized functions such as exponentials\footnote{As outlined in the CUDA programming guide, each streaming multiprocessor (SM) can perform 16 special function operations per clock cycle. The throughput figure arises from multiplying 16 operations by the 132 SMs and an 1830 MHz clock rate, totaling 3.9 TFLOPS for special functions.}—which are critical to softmax calculations. In the FP16 attention forward pass (head dimension 128), the number of FLOPS required for matmul exceeds those needed for exponentiation by a factor of 512; however, the throughput for exponential operations is 256 times less, resulting in the exponential component potentially occupying up to half of the execution time relative to matmul. This imbalance becomes even more pronounced with FP8: the matmul throughput doubles, whereas the throughput for exponentials remains unchanged.

Given that the computation of the exponential function is executed on a distinct hardware resource—the multi-function unit—we ideally seek to overlap this operation with the matmul on the Tensor Cores. To coordinate this, we employ synchronization barriers (\texttt{bar.sync} instructions), ensuring that the GEMM operations (specifically, GEMM1—$\mathbf{P} \mathbf{V}$ from one iteration—and GEMM0—$\mathbf{Q} \mathbf{K}^\top$ from the next) in warpgroup 1 are completed before warpgroup 2 begins its GEMMs. Consequently, while warpgroup 2 is engaged in its GEMM operations, the softmax for warpgroup 1 is computed. The process is then reversed: warpgroup 2 computes softmax while warpgroup 1 returns to performing GEMMs, which is referred to as “pingpong” scheduling. An overview of this scheduling strategy appears in~\cref{fig:pingpong_scheduling}. Although the synchronization in practice may deviate from the conceptual illustration, this approach typically yields improved throughput—such as increasing FP16 forward performance from 570 TFLOPS to approximately 620–640 TFLOPS (with head dimension 128 and sequence length 8192).

\paragraph{Attention variants} In both multi-query attention~\citep{shazeer2019fast} and grouped query attention~\citep{ainslie2023gqa}, our method, as in \textsc{FlashAttention-2}, involves modifying tensor indexing schemes to prevent redundant copies of $\mathbf{K}$ and $\mathbf{V}$ in HBM.

\subsection{Intra-warpgroup overlapping GEMMs and softmax}

Even within one warpgroup, we can overlap some instructions in the softmax with some instructions in the GEMMs. We describe one technique to do so.

Within the attention algorithm, computations inside the main loop exhibit sequential data dependencies, which restrict the extent of intra-iteration parallelism achievable. For instance, the (local) softmax operation (as delineated from lines \ref{code-ws:softmax_start} to \ref{code-ws:softmax_end}) is contingent upon the intermediate outcome $\mathbf{S}_i^{(j)}$ provided by the preceding GEMM operation. Subsequently, the following GEMM receives as input the softmax result, denoted as $\widetilde{\mathbf{P}}_i^{(j)}$. The presence of explicit synchronization points—i.e., wait statements at lines \ref{alg:ws_only_gemm1} and \ref{alg:ws_only_gemm2} in \cref{alg:flash3_wgmma_ws_only}—enforces a serialized execution of the softmax and GEMM routines. To alleviate these constraints, it is possible to introduce a pipelining approach spanning multiple iterations, utilizing supplementary register-based buffers to decouple these operations. Building on this strategy, we introduce a two-stage\footnote{It is important to clarify that while the number of overlapping stages is bounded above by the number $s$ of stages in the circular SMEM buffer, equivalence is not required.} pipelined GEMM-softmax algorithm:

\begin{algorithm}[H]
    \caption{\small\label{alg:flash3_wgmma}\textsc{FlashAttention-3} Consumer Warpgroup Forward Pass}
    \begin{algorithmic}[1]
      \REQUIRE  Input matrices $\mathbf{Q}_i \in \mathbb{R}^{B_r \times d}$ and $\mathbf{K}, \mathbf{V} \in \mathbb{R}^{N \times d}$ residing in HBM, and key block size $B_c$ where $T_c = \lceil \frac{N}{B_c} \rceil$.
      \STATE Adjust register allocation according to the configured number of consumer warps.
      \STATE Initialize on-chip storage for $\mathbf{O}_i = (0) \in \mathbb{R}^{B_r \times d}$ as well as $\ell_i, m_i = (0), (-\infty) \in \mathbb{R}^{B_r}$.
      \STATE Await the loading of $\mathbf{Q}_i$ and $\mathbf{K}_0$ into shared memory.
      \STATE Use WGMMA to calculate $\mathbf{S}_{\mathrm{cur}} = \mathbf{Q}_i \mathbf{K}_0^T$; ensure computation is committed and complete before proceeding.
      \STATE Release the buffer's $0$th stage for $\mathbf{K}$.
      \STATE Derive $m_{i}$, $\tilde{\mathbf{P}}_{\mathrm{cur}}$, and $\ell_{i}$ from $\mathbf{S}_{\mathrm{cur}}$, and appropriately rescale $\mathbf{O}_i$.
      \FOR{$1 \le j < T_c - 1$}
        \STATE Wait until $\mathbf{K}_j$ has been placed into shared memory.
        \label{code:mainloop_start}
        \STATE Perform WGMMA computation: $\mathbf{S}_{\mathrm{next}} = \mathbf{Q}_i \mathbf{K}_{j}^T$, committing the operation without awaiting completion. \label{code:first_wgmma}
        \STATE Wait until $\mathbf{V}_{j-1}$ has been loaded into shared memory.
        \STATE Using WGMMA, update $\mathbf{O}_{i} = \mathbf{O}_{i} + \tilde{\mathbf{P}}_{\mathrm{cur}} \mathbf{V}_{j-1}$; commit this but defer waiting for result. \label{code:second_wgmma}
        \STATE Await completion of the WGMMA calculation $\mathbf{Q}_i \mathbf{K}_{j}^T$.
        \STATE Calculate $m_{i}$, $\tilde{\mathbf{P}}_{\mathrm{next}}$, and $\ell_{i}$ from $\mathbf{S}_{\mathrm{next}}$. \label{code:softmax}
        \STATE After confirming $\tilde{\mathbf{P}}_{\mathrm{cur}} \mathbf{V}_{j-1}$ WGMMA is finished, rescale $\mathbf{O}_i$ accordingly.
        \STATE Release the $(j\,\%\,s)$th and $(j-1\,\%\,s)$th buffer stages for $\mathbf{K}$ and $\mathbf{V}$, respectively.
        \STATE Move $\mathbf{S}_{\mathrm{next}}$ into $\mathbf{S}_{\mathrm{cur}}$.
        \label{code:mainloop_end}
      \ENDFOR \STATE Await loading of $\mathbf{V}_{T_c - 1}$ into shared memory.
      \STATE Compute final update $\mathbf{O}_{i} = \mathbf{O}_{i} + \tilde{\mathbf{P}}_{\mathrm{last}} \mathbf{V}_{T_c - 1}$ via WGMMA, committing and waiting for completion.

\cref{alg:flash3_wgmma} acts as the consumer-side substitute for \cref{alg:flash3_wgmma_ws_only}, together forming the full \textsc{FlashAttention-3} workflow at FP16 precision. In essence, the term WGMMA denotes an asynchronous GEMM in this context. During the main loop (spanning lines \ref{code:mainloop_start} through \ref{code:mainloop_end}), the second invocation of WGMMA within iteration $j$ (see line \ref{code:second_wgmma}) is executed concurrently with the softmax computation belonging to the next iteration, $j+1$ (see line \ref{code:softmax}).

Although the pipelined scheme shown above promises improvements in theoretical throughput, its practical implementation involves several important challenges:

\paragraph{Compiler reordering} Although the pseudocode assumes a specific, ideal sequence of execution, the actual behavior is influenced by the compiler (NVCC), which often rearranges instructions to maximize hardware efficiency. Such rearrangement may interfere with the deliberate interleaving of WGMMA and non-WGMMA operations, potentially causing deviations from the intended pipelining order and thereby undermining anticipated performance benefits. Examination of the SASS output reveals that the compiler does, in practice, produce the desired overlapping pattern (see Section~\ref{sec:2-stage-sass}).

\paragraph{Register pressure} Minimizing register spilling remains crucial for maximizing efficiency. The two-stage pipeline, however, necessitates allocation of additional registers for intermediate values and for maintaining continuity between stages. Notably, it becomes essential to retain an extra $\mathbf{S}_{\mathrm{next}}$ within registers, increasing per-threadblock register use by $B_r \times B_c \times \text{sizeof}(\text{float})$. This heightened demand on registers can compete with the advantages of employing larger block sizes—a widely used method to boost throughput—which itself places high demand on available registers. In practice, choosing between these options requires empirical profiling to determine the optimal trade-off.

\paragraph{3-stage pipelining} Building upon the earlier two-stage pipeline, a three-stage approach is suggested, enabling overlap between the second WGMMA and softmax computation. This can potentially further boost Tensor Core occupancy, but also imposes a greater demand for registers due to the added pipeline stage. As a result, optimizing between increasing tile size and deepening the pipeline becomes even more complex. Detailed discussion and benchmarking of the 3-stage methodology is presented in ~\cref{sec:3-stage}.

\subsection{Low-precision with FP8}
\label{sec:algofp8}

\textbf{Efficiency: layout transformations.} Executing the forward computation of \textsc{FlashAttention-3} using FP8 precision introduces unique obstacles regarding tensor layout that do not arise in FP16 settings.

Typically, the input tensors $\mathbf{Q}$, $\mathbf{K}$, and $\mathbf{V}$ are organized to be contiguous along the head dimension. However, for the second GEMM’s k-major access pattern required by FP8 WGMMA, it is necessary for $\mathbf{V}$—more specifically, the segments of $\mathbf{V}$ accessed within SMEM—to be contiguous along the sequence length axis. Since the TMA load is unable to alter which dimension is contiguous, addressing this issue requires either (1) applying a transpose operation to $\mathbf{V}$ in global memory prior to use, or (2) performing a transpose of the relevant $\mathbf{V}$ tiles within the kernel after these have been loaded into SMEM. To pursue approach (1), possible strategies include (1a) merging the transpose with a prior processing step, such as incorporating it into the rotary embedding’s epilogue, or (1b) invoking a separate kernel specifically to transpose and swap the sequence length and head strides\footnote{A highly tuned transpose kernel can run nearly at device memory bandwidth~\citep{colfax_cutlass_transpose_2024}.}. Nevertheless, (1a) introduces integration complexity into general-purpose libraries, while (1b) is resource-inefficient in inference scenarios that are constrained by memory bandwidth.

In our implementation of FP8 \textsc{FlashAttention-3}, we select approach (2). For the transposition within the kernel, we utilize the LDSM (\verb|ldmatrix|) and STSM (\verb|stmatrix|) instructions, which enable a group of threads within a warp to cooperatively transfer data between shared memory (SMEM) and registers (RMEM) at a 128-byte granularity.\footnote{The PTX specification documents LDSM/STSM as transferring $8 \times 8$ matrices composed of 16-bit elements \cite[\S 9.7.13.4.15-16]{ptx}; by packing two 8-bit elements together, we can also use LDSM/STSM for FP8 data. However, the transpose variants of these instructions are unable to separate packed 8-bit pairs, so additional register shuffling is necessary between the LDSM and STSM stages to accomplish tile-wise transposes; we omit these specifics here.} LDSM and STSM instructions are highly efficient with register usage, making them suitable for execution in the producer warpgroup, and can also alter data layout through the transfer operation itself. Additionally, from the second iteration onward, we can schedule the transpose of the next $\mathbf{V}$ tile such that it runs concurrently with the two WGMMAs computing on the previous $\mathbf{V}$ tile and the current $\mathbf{K}$ tile.

Second, we find that, in contrast to FP16, the FP32 accumulator used in an FP8 WGMMA possesses a register memory layout that does not match that of operand A when stored in registers. This difference is illustrated in \cref{fig:rmem_accum} and \cref{fig:rmem_operand}, where we show how individual register entries are assigned per thread in each layout. Through the application of byte permute instructions, it is possible to reorganize the accumulator produced by the initial WGMMA into a layout that is compatible both for input into a subsequent WGMMA and for alignment with the $\mathbf{V}$ tile layout generated by an in-kernel transpose. In particular, referring to \cref{fig:rmem_accum}, we adjust the register order sequentially to
$$\{ \verb|d0 d1 d4 d5 d2 d3 d6 d7| \},$$
repeating this permutation for every consecutive group of 8 bytes. At the level of the $\mathbf{P}$ tile's logical structure, this operation rearranges its columns (e.g., columns $0189$ shift to become the leading set of columns). To ensure the next WGMMA receives the correct data arrangement, the in-kernel transpose can likewise output a reordered set of rows for the $\mathbf{V}$ tile. This added flexibility from the in-kernel transpose removes the necessity for inter-thread shuffling to redistribute registers, as was required in our earlier approach~\citep{colfax_fp8_flashattention_2024}.

\textbf{Accuracy: block quantization and incoherent processing.} The FP8 (e4m3) representation allocates only 3 bits for the mantissa and 4 for the exponent, which increases the susceptibility to numerical inaccuracies compared to FP16 or BF16. In addition, large-scale models frequently produce outlier activations~\citep{dettmers2208llm, sun2024massive}—values whose magnitudes vastly exceed the majority—posing a challenge for quantization. Standard practice involves per-tensor scaling~\citep{micikevicius2022fp8}, wherein each tensor (such as $\mathbf{Q}$, $\mathbf{K}$, or $\mathbf{V}$) is associated with its own scaling factor. To improve the numerical robustness of attention operations under FP8 quantization, we introduce two specific strategies:

\begin{enumerate}

\item \textbf{Block quantization}: In this approach, a single scalar is retained per block; specifically, for each of $\mathbf{Q}$, $\mathbf{K}$, and $\mathbf{V}$, the tensor is divided into blocks of dimensions $B_r \times d$ or $B_c \times d$, with quantization performed independently for each block. This process can be integrated with the operation that immediately precedes attention computation (such as rotary embedding), incurring no extra latency because the rotary embedding step is limited by memory bandwidth. Given that the \textsc{FlashAttention-3} algorithm natively operates on block structures, each corresponding block of
  $\mathbf{S}$ can be rescaled to accommodate the effects of block quantization, and this can be achieved without any additional computational overhead.

\item \textbf{Incoherent processing}: To mitigate the influence of outlier values, $\mathbf{Q}$ and $\mathbf{K}$ are pre-multiplied by a random orthogonal matrix $\mathbf{M}$ prior to quantization into FP8 format. Since $\mathbf{M}$ is orthogonal, it satisfies $\mathbf{M} \mathbf{M}^\top = I$, and therefore $(\mathbf{Q} \mathbf{M})(\mathbf{K} \mathbf{M})^\top = \mathbf{Q}\mathbf{K}^\top$, ensuring that the attention result is unchanged by this transformation. The effect of this technique is to distribute outliers, as each element of $\mathbf{Q}\mathbf{M}$ or $\mathbf{K}\mathbf{M}$ consists of a random linear combination of the elements of $\mathbf{Q}$ or $\mathbf{K}$, leading to a decrease in quantization errors. Following the implementations in \citet{chee2024quip} and~\citet{tseng2024quip}, we select $\mathbf{M}$ as a product of randomly signed diagonal matrices and a Hadamard matrix, which enables efficient multiplication in $O(d \log d)$ time as opposed to the more costly $O(d^2)$, and allows this process to be fused with rotary embedding with no additional computation required.
\end{enumerate}

We validate that these two techniques reduces numerical error by up to 2.6$\times$ in \cref{sec:numerical_error}.

\section{Empirical Validation}
\label{sec:exp}

We use the primitives from CUTLASS~\citep{Thakkar_CUTLASS_2023} such as WGMMA and TMA abstractions to implement \textsc{FlashAttention-3} and evaluate its efficiency and accuracy.

\begin{itemize}

\item \textbf{Benchmarking attention.} We evaluate the performance of \textsc{FlashAttention-3} in terms of runtime over various sequence lengths and conduct comparisons with several alternatives: a conventional PyTorch implementation, \textsc{FlashAttention-2}, a Triton-based variant of \textsc{FlashAttention-2} (which leverages H100-tailored instructions), and the cuDNN-optimized vendor implementation of \textsc{FlashAttention-2} for H100 GPUs. Our results demonstrate that \textsc{FlashAttention-3} can achieve speedups of up to 2.0$\times$ relative to \textsc{FlashAttention-2} and up to 1.5$\times$ in comparison with the Triton implementation. Moreover, \textsc{FlashAttention-3} attains up to 740 TFLOPs/s, corresponding to 75\% of the H100 GPU's peak theoretical TFLOPs/s.
\item \textbf{Ablation study.} Through further experiments, we establish that the performance gains in \textsc{FlashAttention-3} can be attributed to advancements such as warp-specialization and the integration of GEMM-softmax pipelining within our algorithmic design.
\item \textbf{Accuracy of FP8 attention.} Our findings also indicate that employing block quantization along with incoherent processing significantly mitigates numerical errors, achieving a 2.6$\times$ reduction in FP8 \textsc{FlashAttention-3}.

\subsection{Benchmarking Attention}
\label{subsec:benchmark_attn}

We evaluate the execution times of various attention algorithms on an H100 80GB SXM5 GPU across multiple configurations, specifically toggling the causal mask and varying the head dimension between 64 and 128, utilizing FP16 precision for inputs. The performance outcomes, presented in~\cref{fig:benchmark_attn_fwd} and~\cref{fig:benchmark_attn_bwd}, reveal that \textsc{FlashAttention-3} achieves a speedup of roughly 1.5 to 2.0$\times$ over \textsc{FlashAttention-2} in the forward computation, and accelerates the backward computation by a factor of 1.5 to 1.75$\times$. Relative to conventional attention implementations, \textsc{FlashAttention-3} exhibits speed improvements in the range of 3$\times$ to 16$\times$. Notably, when working with sequence lengths of 1,000 tokens or greater, \textsc{FlashAttention-3} outperforms the closed-source vendor library cuDNN, which is tailored for the H100 GPU architecture.

\paragraph{Benchmark settings:} Sequence lengths are set to 512, 1k, up to 16k, while the batch size is adjusted such that the aggregate token count remains at 16k. The hidden dimension is fixed at 2048, and the head dimension is varied among 64, 128, and 256, corresponding to 32, 16, or 8 attention heads, respectively. The FLOPs for the forward pass are computed as follows:

\begin{equation*}
  4 \cdot \text{seqlen}^2 \cdot \text{head dimension} \cdot \text{number of heads}.
\end{equation*}

When applying causal masking, the total count is halved, as roughly 50\% of the possible entries are computed. To estimate the backward pass FLOPs, we scale the forward pass FLOPs by a factor of 2.5, reflecting that the forward pass involves two matrix multiplications, while the backward pass entails five, owing to recomputation.

Additionally, we assess the forward pass runtime in FP8 using analogous experimental conditions. Performance for a headdim of 256 is presented in~\cref{fig:benchmark_attn_fp8}, with comprehensive outcomes provided in \cref{sec:benchmark_attn_fp8_full}.

\subsection{Ablation Study: 2-Stage Pipelining Experiments}
\label{sec:2-stage-pipelining-experiments}

We conduct an ablation study on both the 2-stage WGMMA-softmax pipeline and warp-specialization in the context of non-causal FP16 \textsc{FlashAttention-3}, using constant settings for $\{\text{batch}, \text{seqlen}, \text{nheads}, \text{hdim}\} = \{ 4, 8448, 16, 128\}$. As presented in~\cref{table:ablation_pipelining}, these algorithmic enhancements—namely, introducing asynchrony through warp-specialization and enabling concurrent execution of GEMM and softmax operations—produce a notable performance gain, elevating throughput from 570 to 661 TFLOPs.

\subsection{Numerical Error Validation}
\label{sec:numerical_error}

Given ongoing concern regarding the numerical errors associated with \textsc{FlashAttention}~\citep{golden2024flash}, we carry out a comparison among \textsc{FlashAttention-2}, \textsc{FlashAttention-3}, and a baseline attention implementation, all evaluated against a high-precision FP64 reference. To mimic the outlier features and activations frequently observed in large language models (LLMs)~\citep{dettmers2208llm, sun2024massive}, we sample the elements of $\mathbf{Q}, \mathbf{K}, \mathbf{V}$ using the following probability distribution:

\begin{equation*}
  \mathcal{N}(0, 1) + \mathcal{N}(0, 100) \cdot \mathrm{Bernoulli}(0.001).
\end{equation*}

Specifically, all entries are sampled from a normal distribution with mean zero and standard deviation 1; however, for 0.1\% of the entries, an independent normal noise with standard deviation 10 is added. The root mean squared error (RMSE) is reported in~\cref{table:numerical_error}. When using FP16, both \textsc{FlashAttention-2} and \textsc{FlashAttention-3} yield an RMSE that is 1.7$\times$ lower than that of the conventional approach, as intermediate softmax values are stored in FP32. For FP8, the baseline attention method uses per-tensor scaling, computes matmuls with an FP32 accumulator, and maintains softmax intermediates in FP16. Owing to its block quantization and incoherent processing strategies, \textsc{FlashAttention-3} with FP8 achieves 2.6$\times$ greater accuracy compared to this baseline.

\section{Dicussion, Limitations, Conclusion}
\label{sec:discussion}

In developing \textsc{FlashAttention-3}, we have shown that innovative programming strategies and leveraging hardware advancements—particularly asynchronous operations and low-precision computation—can substantially boost both the efficiency and precision of attention mechanisms. Our approach achieves a 1.5-2.0$\times$ speedup over \textsc{FlashAttention-2} and offers a 2.6$\times$ reduction in FP8 numerical error relative to conventional per-tensor quantization approaches. Looking ahead, there are several areas for improvement: enhancing optimization specifically for LLM inference workloads, adopting a persistent kernel design for the FP8 kernel,\footnote{For benchmarking purposes, FP16 \textsc{FlashAttention-3} benefits from a persistent kernel along with a load-balancing approach, whereas FP8 \textsc{FlashAttention-3} currently lacks these features. This is a contributing factor to the lower performance of FP8 \textsc{FlashAttention-3} at shorter sequence lengths and with causal masking, compared to FP8 cuDNN kernels.} and investigating the implications of employing low-precision attention in large-scale training regimes. While our study centers on Hopper GPUs, the concepts and methods introduced here should be adaptable to other accelerator hardware. Ultimately, we anticipate that improved attention primitives—delivering both higher speed and accuracy—will pave the way for advancements in applications involving extended context lengths.

\end{document}